%%*************************************************************************
%%
%% Křívky učení v raném stadia učení, anebo setkání člověka s cizím nastrojem.
%% V0.0
%% 2012/11/28
%% Timothy Hobbs
%% timothyhobbs@seznam.cz
%%
%%*************************************************************************
%%
%% Legal Notice:
%%
%% This code is offered as-is without any warranty either expressed or
%% implied; without even the implied warranty of MERCHANTABILITY or
%% FITNESS FOR A PARTICULAR PURPOSE! 
%% User assumes all risk.
%%
%% This work by Timothy Hobbs is licensed under a
%% Creative Commons Attribution-NonCommercial-ShareAlike 3.0 Unported License.
%% http://creativecommons.org/licenses/by-nc-sa/3.0/
%%
%% Loriem ipsum was from
%% https://github.com/pborky/documents/tree/master/DIP
%% by Peter Boraros
%%*************************************************************************
\documentclass[a4paper,twoside,english,english,openright,cleardoublepage=empty,BCOR10mm,DIV11]{scrreprt}
\usepackage[T1]{fontenc}
\usepackage[utf8]{inputenc}

\usepackage[czech]{babel}%počeštění názvů (Obsah, Kapitola, Literatura atp.)
%<------------------------------ headers ------------------------------------>
\newcommand{\SubjectCs}{Bakala\v{r}sk\'{a} Prace}
\newcommand{\TitleCs}{K\v{r}\'{i}vky u\v{c}en\'{i} v ran\'{e}m stadia u\v{c}en\'{i}, anebo setk\'{a}n\'{i} \v{c}lov\v{e}ka s ciz\'{i}m nastrojem.}
\newcommand{\TitleEn}{Learning curves in early stages of skill aquisition.}
\newcommand{\KeywordsEn}{Learning Curves, Braille, Reading Speed}
\newcommand{\KeywordsCs}{K\v{r}\'{i}vky u\v{c}en\'{i}, Braillové pismo, rychlost čtení}
\newcommand{\AuthorName}{Timothy Vladim\'{i}r Hobbs}
\newcommand{\SupervisorName}{Ing. Karolina Duschinsk\'{a}, Ph.D.}
\newcommand{\SupervisorEmail}{karolina.duschinska@pedf.cuni.cz}
\newcommand{\IssuedIn}{Prague}
\newcommand{\DepartmentEn}{Department of Pedagogy}
\newcommand{\DepartmentCs}{Katedra pedagogiky}
\newcommand{\Thanks}{Na tomto m\'{i}st\v{e} mohou b\'{y}t naps\'{a}na p\v{r}\'{i}padn\'{a} pod\v{e}kov\'{a}n\'{i} %%
(vedouc\'{i}mu pr\'{a}ce, konzultantovi, tomu kdo p\r{u}j\v{c}il software, literaturu, %%
poskytl data apod.).}

%<------------------------------         ------------------------------------>

\usepackage{array}
\usepackage{longtable}
\usepackage{varioref}
\usepackage{wrapfig}
\usepackage{fancybox}
%\usepackage{calc}
\usepackage{framed}
\usepackage{url}
\usepackage{amssymb,amsmath}

%\usepackage{graphicx}

\usepackage{ifpdf}
\ifpdf
\usepackage[pdftex]{graphicx}
\graphicspath{{./img/}}
\DeclareGraphicsExtensions{.pdf}
\else
\usepackage[dvips]{graphicx}
\graphicspath{{./img/}}
\DeclareGraphicsExtensions{.eps}
\fi

\makeatletter

%%%%%%%%%%%%%%%%%%%%%%%%%%%%%% LyX specific LaTeX commands.
\providecommand{\LyX}{L\kern-.1667em\lower.25em\hbox{Y}\kern-.125emX\@}
\newcommand{\lyxline}[1][1pt]{%
  \par\noindent%
  \rule[.5ex]{\linewidth}{#1}\par}
\newcommand{\noun}[1]{\textsc{#1}}
%% Special footnote code from the package 'stblftnt.sty'
%% Author: Robin Fairbairns -- Last revised Dec 13 1996
\let\SF@@footnote\footnote
\def\footnote{\ifx\protect\@typeset@protect
    \expandafter\SF@@footnote
  \else
    \expandafter\SF@gobble@opt
  \fi
}
\expandafter\def\csname SF@gobble@opt \endcsname{\@ifnextchar[%]
  \SF@gobble@twobracket
  \@gobble
}
\edef\SF@gobble@opt{\noexpand\protect
  \expandafter\noexpand\csname SF@gobble@opt \endcsname}
\def\SF@gobble@twobracket[#1]#2{}
%% Because html converters don't know tabularnewline
\providecommand{\tabularnewline}{\\}

%%%%%%%%%%%%%%%%%%%%%%%%%%%%%% Textclass specific LaTeX commands.
\newenvironment{lyxcode}
{\begin{list}{}{
\setlength{\rightmargin}{\leftmargin}
\setlength{\listparindent}{0pt}% needed for AMS classes
\raggedright
\setlength{\itemsep}{0pt}
\setlength{\parsep}{0pt}
\normalfont\ttfamily}%
 \item[]}
{\end{list}}

%%%%%%%%%%%%%%%%%%%%%%%%%%%%%% User specified LaTeX commands.
%<-------------------------------společná nastavení------------------------------>
\usepackage[]{hyperref} %odkazy v  pdf jsou klikací s barevnými rámečky
%\usepackage[urlcolor=blue]{hyperref}
%\usepackage[colorlinks=true,urlcolor=blue]{hyperref}
\usepackage[numbers,sort&compress]{natbib} %balíček pro citace literatury  
\usepackage{hypernat}%interakce mezi hyperref a natbib
\newcommand{\BibTeX}{{\sc Bib}\TeX}%BibTeX logo
\hypersetup{   % Nastavení polí PDF dokumentu 
pdftitle={\TitleCs},%   
pdfauthor={\AuthorName},%  
pdfsubject={\SubjectCs},%   
pdfkeywords={\KeywordsCs}%                             
}
\usepackage{multicol}

%<-----------------------------volání stylů----------------------------------------->
% (znak % je označení komentáře: co je za ním, není aktivní)
%<------------------------------------písmo----------------------------------------->
%\usepackage{packages/bc-latinmodern}
%\usepackage{packages/bc-times}
\usepackage{packages/bc-palatino}
%\usepackage{packages/bc-iwona}
%\usepackage{packages/bc-helvetika}


%<------------------------------záhlaví stránek------------------------------------>
%\usepackage{packages/bc-headings}
\usepackage{packages/bc-fancyhdr}

%<------------------------------hlavičky kapitol------------------------------------>
%\usepackage{packages/bc-neueskapitel}
\usepackage{packages/bc-fancychap}

\makeatother

\usepackage{babel}



\begin{document}
~\thispagestyle{empty}{\small ~\vfill{}
}{\small \par}

\noindent {\small \Thanks \newpage{}}{\small \par}

~\thispagestyle{empty}\vfill{}
%Tato stránka je tzv. protititul a je graficky součástí titulní stránky.
%Nechte ji prázdnou, nebo na ni umístěte vhodnou fotografii či ilustraci.

\cleardoublepage{}~\thispagestyle{empty}\begin{center}\pagenumbering{roman}\vspace{10mm}


\textsf{\textsc{\noun{\LARGE Universita Karlova v Praze}}}\\
\vspace{0.5em}
\textsf{\textsc{\noun{\LARGE Pedagogická Fakulta}}}\\
\vspace*{1em}
\textsf{\textsc{\noun{\Large \DepartmentCs}}}\vspace{15mm}


\includegraphics[width=0.3\textwidth]{uk}\vspace{15mm}

\textsf{\textsc{\noun{\huge \SubjectCs}}}{\huge \par}

\vspace{15mm}


\textsf{\textsc{\noun{\LARGE \TitleCs}}}{\LARGE \par}


\vspace{10mm}


\end{center} 

\vspace*{\fill}


\vspace{10mm}


\begin{description}
\item [{{\large Autor:}}] \noindent \textsf{\large \AuthorName}{\large \par}
\item [{{\large Vedoucí}}] \noindent \textsf{\large \SupervisorName}%
{\large \hfill{}}\textsf{\large \IssuedIn, \the\year}{\large{}
}{\large \par}
\end{description}
\clearpage{}

{\small \thispagestyle{plain}\addcontentsline{toc}{chapter}{Abstract} }{\small \par}

\noindent {\small ~\vfill{}
}{\small \par}

\begin{description}
\item [{{\small Název~práce:}}] \noindent {\small \TitleCs}{\small \par}
\item [{{\small Autor:}}] \noindent {\small \AuthorName}{\small \par}
\item [{{\small Katedra~(ústav):}}] \noindent {\small \DepartmentCs}{\small \par}
\item [{{\small Vedoucí~diplomové~práce:}}] \noindent {\small \SupervisorName}
\item [{{\small e-mail~vedoucího:}}] \noindent {\small \SupervisorEmail}\\
{\small \par}
\item [{{\small Abstrakt}}] \noindent {\small Zkoumání prvních dojmů, adaptace, a učení během kratké doby a s neznámým materiálem.
V teto studii, jsem zkoumal jak reagují na nový druh braillsého řádku FCHAD a jak se ho naučí používat. Pedagogické prostředí teto studie je přísné a učení je založené na drilování. Zkoumání se zvlašťě zaměří na otázku adaptace během přísného drilování.  Studie se podobá psychologickým výzkumům křivek učení s tím rozdílem, že pedagog je plnohodnotná součást studie.
 }{\small \par}
\item [{{\small Klíčová~slova:}}] \noindent {\small \KeywordsCs}\\
{\small \lyxline{\small}}{\small \par}
\item [{{\small Title:}}] \noindent {\small \TitleEn}{\small \par}
\item [{{\small Author:}}] \noindent {\small \AuthorName}{\small \par}
\item [{{\small Department:}}] \noindent {\small \DepartmentEn}{\small \par}
\item [{{\small Supervisor:}}] \noindent {\small \SupervisorName}{\small \par}
\item [{{\small Supervisor's~e-mail~address:}}] \noindent {\small \SupervisorEmail}\\
{\small \par}
\item [{{\small Abstract}}] \noindent {\small An examination of first impressions, adaptation and learning over a short period with novel material.
 In this study, I have examined how subjects react to, and learn to use, an FCHAD braille display.  The pedagogic atmosphere of this study is strict and drill based.  Particular focus is placed on how students adapt during the course of a drill.  The study was performed similarly to many learning curve studies, yet differs in that a teaching element was explicitly included rather than explicitly excluded.
 }{\small \par}
\item [{{\small Keywords:}}] \noindent {\small \KeywordsEn}{\small \par}
\end{description}
\cleardoublepage{}\thispagestyle{empty}~{\small \addcontentsline{toc}{chapter}{Assignment of the Work} }{\small \par}



\newpage{}\thispagestyle{empty}~


\newpage{}\thispagestyle{plain}

{\small %\setcounter{page}{3} % nastavení číslování stránek
\ }{\small \par}

\noindent {\small \vfill{}
 % nastavuje dynamické umístění následujícího textu do spodní části stránky
~}{\small \par}

\noindent {\small Prohlašuji, že jsem svou diplomovou práci napsal(a)
samostatně a výhradně s použitím citovaných pramenů. Souhlasím se
zapůjčováním práce a jejím zveřejňováním.}{\small \par}

{\small \bigskip{}
}\noindent {\small{} V Praze dne \today\hspace{\fill} \AuthorName}\\
{\small{} % doplňte patřičné datum, jméno a příjmení
}{\small \par}

{\small %%%   Výtisk pak na tomto míste nezapomeňte PODEPSAT!
%%%                                         *********
}{\small \par}

\cleardoublepage{}\thispagestyle{empty}{\small \tableofcontents{}% vkládá automaticky generovaný obsah dokumentu
\cleardoublepage{}}{\small \par}


\chapter{Úvod}

\section{Nevidomost na světě a význam Braillova písma}

Nevidomost je fenomen rozvojových zemí a starších lidech. 39 milionů nevidomých lidí žijí na světě. Devadesát procent z nich bydli v rozvojových zemích a 65\% jsou starší 50 let\citep{whodata}.

Gramotnost nevidomých je schopnost čist Braillova písma.  Mezi sebou nevidomé říkají číst rukama.

V vyspělých zemí gramotnost souvislí s zaměstnaností. Podle statistiky zveřejněné v spojených státech 90\% procent gramotných nevidomých lidí jsou zaměstnání oproti jenom 30\% negramotných nevidomých lidí.

Nenašel jsem podobná čísla pro české republiky ani rozvojový svět.  Do jisté míru statistiky jsou zkreslené.  V státech gramotnost značně poklesá.  Padesát procent nevidomých amerických žáků se učili číst rukama v roce 1960. Dneska jenom dvanáct procent.  Dříve mnohou lidí ztratili zrak kvůli infekčním nemoci dneska skoro nikdo. Dále už umíme léčit bílý zákal, a ve vyspělých zemí ne-existuje nevidomost kvůli bílých zákalům.  Jsou miň nevidomé lidi ale ty nevidomé lidi, kteří jsou jsou vice nemocné.  Daleko větší proporce nevidomých lidí ve vyspělém světě mají vedlejší nemoc, která překáží gramotnosti.  Mnohou nevidomých trpěli předčasného porodu anebo jiné perinatalním závad\citep{perkins,whozakal,whodata}.

Je další důvod proč gramotnost nevidomých je tak nízká.  Materiály v braillském písmu jsou velmi omezený.  V Knihovně a tiskárna pro nevidomé K.E.Macana jsou 5693 titulů tisknutí braillském písmu\citep{biblio}. V Městské Knihovně v Praze jsou 387625 knih pro vidící \citep{mlp}.

Složení nevidomost v rozvojových zemí je výrazně jiná než ve vyspělém světě. Osmdesát procent nevidomost na světě je či léčitelná anebo nevyhnutelná\citep{whodata}.  Však celý 50\% nevidomosti na světě je kvůli bílému zákalu\citep{whozakal}.

Bílý zákal není jen léčitelný. Je to docela snadno a levně léčitelný. Chirurgie pro bílý zákal soji mezi 1 000 a 1 700 USD v Indii\citep{cataractsindia}. Z toho víme, že většina nevidomých lidí na světě jsou velmi chudé.
\clearpage


\section{Co to je Braillovo písmo?}

Braillovo písmo je hmatní abeceda, které lze číst a psát bez zraku.

Každé písmeno má dvě sloupce zvýšených bodů. Podle toho kolik bodů jsou zvýšené a kde jsou mezery mezi body víme, které písmeno je napsáno.

Braillovo písmo je jiné v každém země.  V Čechách používáme šestibodové písmo pro tisk a osmibodové písmo na počítače.  V Německu písmo je vždy osmibodové.

Abychom mohli mluvit o braillských písmenech existuje číselní označení pro body.


\begin{table}[!hb]
\centering
\makebox[0pt][c]{\parbox{1\textwidth}{%
\begin{minipage}[b]{0.3\hsize}\centering
% putStrLn $ utf8ToBrailleBox "1⠁2⠂3⠄4⠈5⠐6⠠7⡀8⢀"
\begin{tabular}{|c|c|}
\hline
1\braillebox{1       }&4\braillebox{   4    }\\
2\braillebox{ 2      }&5\braillebox{    5   }\\
3\braillebox{  3     }&6\braillebox{     6  }\\
7\braillebox{      7 }&8\braillebox{       8}\\
\hline
\end{tabular}

číselní označení bodů

\end{minipage}
\hfill
\begin{minipage}[b]{0.68\hsize}\centering

%%%%%%%%%%%


% putStrLn $ utf8ToBrailleBox "a ⠁ b ⠃ c ⠉ d ⠙ e ⠑ f ⠋ g ⠛ h ⠓ i ⠊ j ⠚ k ⠅ l ⠇ m ⠍ n ⠝ o ⠕ p ⠏ q ⠟ r ⠗ s ⠎ t ⠞ u ⠥ v ⠧ x ⠭ y ⠽ z ⠵ ý ⠯ w ⠷ ž ⠮ ů ⠾ á ⠡ ě ⠣ č ⠩ ď ⠹ š ⠳ ň ⠫  ť ⠳ ó ⠪ ř ⠺ í ⠌ é ⠜ ú ⠬ "
\begin{tabular}{|c|c|c|c|c|c|c|c|}
\hline
a \braillebox{1       }&b \braillebox{12      }&c \braillebox{1  4    }&d \braillebox{1  45   }&e \braillebox{1   5   }&f \braillebox{12 4    }&g \braillebox{12 45   }&h \braillebox{12  5   }\\
\hline
i \braillebox{ 2 4    }&j \braillebox{ 2 45   }&k \braillebox{1 3     }&l \braillebox{123     }&m \braillebox{1 34    }&n \braillebox{1 345   }&o \braillebox{1 3 5   }&p \braillebox{1234    }\\
\hline
q \braillebox{12345   }&r \braillebox{123 5   }&s \braillebox{ 234    }&t \braillebox{ 2345   }&u \braillebox{1 3  6  }&v \braillebox{123  6  }&x \braillebox{1 34 6  }&y \braillebox{1 3456  }\\
\hline
z \braillebox{1 3 56  }&ý \braillebox{1234 6  }&w \braillebox{123 56  }&ž \braillebox{ 234 6  }&ů \braillebox{ 23456  }&á \braillebox{1    6  }&ě \braillebox{12   6  }&č \braillebox{1  4 6  }\\
\hline
ď \braillebox{1  456  }&š \braillebox{12  56  }&ň \braillebox{12 4 6  }&ť \braillebox{12  56  }&ó \braillebox{ 2 4 6  }&ř \braillebox{ 2 456  }&í \braillebox{  34    }&é \braillebox{  345   }\\
\hline
ú \braillebox{  34 6  }\\
\hline
\end{tabular}
% -- https://github.com/timthelion/utf8-to-braillebo://github.com/timthelion/utf8-to-braillebox

česká braillská abeceda

\end{minipage}%
}}
\end{table}


Dejte pozor na to, že v abecedě nejsou velké písmo ani čísla. Ve šestibodové Braillova písma velká písmena jsou escapování. Písmenech píšeme velký A jako
% putStrLn $ utf8ToBrailleBox "⠠⠁"
\braillebox{     6  }\braillebox{1       }
%%%%
. Samotní šestí bod je escape sekvence pro velká písmena. Čísla jsou taky escapování. Escape sekvence pro čísla je
% putStrLn $ utf8ToBrailleBox "⠼"
\braillebox{  3456  }
%%%%
tak, že čísla v šestibodové braillských písmenech jsou:

% putStrLn $ utf8ToBrailleBox "1⠼⠁2⠼⠃3⠼⠉4⠼⠙5⠼⠑6⠼⠋7⠼⠛8⠼⠓9⠼⠊0⠼⠚"
1\braillebox{  3456  }\braillebox{1       }
2\braillebox{  3456  }\braillebox{12      }
3\braillebox{  3456  }\braillebox{1  4    }
5\braillebox{  3456  }\braillebox{1   5   }
6\braillebox{  3456  }\braillebox{12 4    }
7\braillebox{  3456  }\braillebox{12 45   }
8\braillebox{  3456  }\braillebox{12  5   }
9\braillebox{  3456  }\braillebox{ 2 4    }
0\braillebox{  3456  }\braillebox{ 2 45   }
%%%%

Osmibodové Braillovo písmo nemá escape sekvence.  Sedmi a osmi body se používají k číslům a kapitalizaci. Velký A je
% putStrLn $ utf8ToBrailleBox "⡁"
\braillebox{1     7 }
%%%%
. A čísla jsou podznačené osmým bodem: 1 je
% putStrLn $ utf8ToBrailleBox "⢁"
\braillebox{1      8}
%%%%
. Ne-existuje žádný standard pro osmibodové braillské písmo\citep{6dotbraille}.  Protože osmibodové braillské písmo je obecně zobrazené na počítače uživatele mohou nastavit vlastní kódování.

Existuje Braillovo kódování pro hudební noty i matematické rovnice.

\section{Jak funguje braillský řádek a jak funguje FCHAD?}

\subsection{Braillský Řádek}

Už jsem zmínil, že existuje málo braillských knih. Braillovo písmo má podstatu hlavně jako grafické rozehraní u počítačů.  Mnoho nevidomých lidí používají hlasový vystup ale aby mohly pochopit a upravit složité formátované texty potřebují braillský řádek.

Braillský řádek je elektromechanické zařízení, který \uv{zobrazuje} jeden řádek braillského textu.  Braillský řádek stojí 2 až 8 tisíc dolarů a zobrazí 40 až 80 znaků.  Obecně cena stoupá s výším počtem znaků.  Braillský řádek, který zobrazuje jen 12 znaků stoji jen tisíc dolarů\citep{perkinsdisplays}.

Dnešní braillské řádky používají piezoelektrickou technologie.  Piezoelektrické krystaly jsou krystaly, které změní velikost když jím projde elektrický proud. Piezoelektrické řádky mají 16 dlouhých krystalů za každý znak.  Samotné krystaly jsou dost drahé a to je vidět v ceně braillského řádku ale to není jediný problém.  Ještě větší cena než cena samotného řádku je cena údržby.  Prach, který vleze do malé díry v řádce muže poškodit piezoelektrický mechanismus.  Když kupujete řádek od americké nezískovky Perkins například byste měl taky zařídit smlouvu na údržbu, která stoji mezi 300 a 500 dolarů USD\citep{perkinsdisplays}.

\subsection{FCHAD}

Už par let pracují na nový vynález tzn. FCHAD. FCHAD je jednoduší varianta braillského řádku, který je odolnější a levnější.  FCHAD, který jsem používal v studie stalo miň než sto dolarů vyrobit.

Není to žádná genialita, že jsem dokázal vyrobit řádek, který stoji o tolik krát miň.  FCHAD řádek je jednoznačně horší než peizoelektrický řádek.  Čte se na to pomalejší a je nutný používat oba ruce k čtení.  Hlavně, že je to levný!

Najdete obrázky FCHADu v apendixu A2 kde jsou zobrazené dvě verzi FCHADu, které byly použity v studie.

FCHAD je zkratka pro \uv{Fast Character Display} tj rychlý zobrazovač písmen.  FCHAD umí zobrazit jenom jedno písmeno naráz ale to neznamená, že nemůžete přečíst delší texty s FCHADem.  FCHAD má ještě další komponent a to je selektor.  Obrázek 1 B a C jsou dvě různé selektory.  Selektor může taky byt počítačový myš.  Selektorem můžete změnit písmeno, které je zobrazené.

Už někomu napadlo, že by to bylo levný zobrazit jen jedno písmeno. Však mnoho FCHADu již existují ale žádné na trhu nejsou.  Můžete prohlídnout různé FCHADy v sekce \uv{Analýza a srovnání.}

\section{Význam pedagogiky v oblasti zařízení pro nevidomé}

Uznávám, že nevím proč tolik FCHADu byly navržené aniž by dostali na trhu.  Když jsem se zeptal různé lidí, které s Braillským písmenem pracovali jsem slyšel často cynický názor, že je větší zisk vyrobit drahý piezo než levný FCHAD.

Mnoho vynálezců pracují v rámci akademie. Sám nemají možnost založit firmu. Třeba vynalezli ty nastroje ale nenašli nikoho, který je chtěl vyrobit. Musím říct, že je třeba hodně trpělivost a stálost v práci přenášet nový vynález na trhu.  Nejsem ani polovinu cestu a už mě trvalo vice než rok.

Velázquez, kdy píše o přenosných pomocníků nevidomých má jiný názor proč tolik zařízení se ztratí na cestě.  V závěru vlastní studie píše:
\em
\uv{Přes obrovské úsilí badatelů, a přes obrovské množství technologii na trhu, nevidomých lidí odmítají novou technologie. Audio knihy, braillské řádky, bílý hůl, a vodicí pes jsou stále nejoblíbenější technologii u nevidomých.

Přijatý přenositelných technologie je těžký úkolů pro badatele.  Motivace, kooperace, optimismus, ochotnost/schopnost se učit anebo přizpůsobit se není nikdy samozřejmost u nevidomých.}\em \citep[str. 15, přeložený z angičtiny]{velazquez2010wearable}%Despite efforts and the great variety of wearable assistive devices available, user acceptance is quite low. Audio books and Braille displays (for those who can read Braille) and the white cane and guide dog will continue to be the most popular reading/travel assistive devices for the blind.
% Acceptance of any other portable or wearable assistive device is always a challenge in blind population. Motivation, cooperation, optimism, willingness/ability to learn or adapt new skills is not a combination that can be taken for granted.

Sám jsem zažil neochota od nevidomých lidí. První nevidomý člověk, který jsem ukazoval stroj vůbec nepochopil jak to má fungovat ale byl jistý, že to není dobrý.

Spousta technologie musejí čekat až lidi jsou na ně zvykle aby mohly byt použité.  Pamatují si na mámu a jak ona stále hledá na mapě na papírové mapě cestu, kterou by našla několik krát rychlejší přes internetem.  A myslím o tom jak dlouho ještě hledala v telefonním seznamu a pak volala pro informace, kterou lze jednoduše přečíst na internetu.  Vidím jasně, že zásadnější část než všimneme, technologický rozvoj je řízený technologické učení a zvyklost světové obyvatelství.  Pedagog a pedagogické teorii sobě najdou uprostřed velké civilizační pokroky a i když pedagogové samotné jsou často nejhorší na tom, když to jde o adaptace na novou technologii, učit a učit i přes neochota žáka je úkolem pedagoga, a je úkolem pedagoga pomáhat žákovi navigovat světě, které je ve stavu trvalé změny.  Ještě štěstí většina nevidomých lidí ochotní se učit jsou.

\section{Co vlastně vyzkoušíme?}

Moje otázka je, jak lidí reagují na neznámý nastroj během první 45 minut použity?  Jak oni přizpůsobují k novému nastroje?  Jaké mají potom dojmy?  Lze z těchto setkání taky kreslit křivky učení.

Zajímavý vlastnost studie je v tom, že FCHAD je úplně neznámý pro účastnici.  Nikdy v životě neslyšeli zkratu FCHAD a nikdy nesetkali dříve s podobným nástroji.  Použity FCHADu je složitý motorický a konceptuálně.

Seděl jsem s šesti nevidomým lidmi vyzkoušet svůj FCHAD. Účastníci četli krátký text s FCHADem a pak řekli své dojmy.  První čtyři krát jsem dorazil na technické potíže, ale mám platný výsledky z tři účastníků.



\chapter{Anomaly Detection}

%% this is intended to be brief, so keep it simple and general
In general an anomaly detection is the problem of finding patterns in data that do not 
conform to expected behavior.
A term \emph{anomaly} or \emph{outlier} refers to these non-conforming patterns. 
Usually knowledge about non-conforming patterns is inportant due to fact that they may refer
to singnificant information, in many cases also critical and actionable, 
e.g. a tumor presence may be indicated by anomalous magnetic resonance imaging (MRI) scan, 
network intrusion may cause anomalous signature of the packets to be observed.

%% history
The anomaly detection has been studied as early as the $19^{th}$ century by 
statisticians as a statistical method.
Due now, several techniques have been developed, using domain-independent approach 
or developed specificaly for particular domain. %TODO moar faking history

Apparently simple approach of anomaly detection is to define a region representing 
normal behavior and declare any patterns which does not conform to this region as anomaly. 
This approach is obfuscated by several factors:
\begin{itemize}
	\item Definition of normal behavior must contain every possible normal behavior 
	and it is difficultly achievable.
	\item The boundary between anomalies and normal behavior is not accurate and 
	can introduce wrong interpretation	of particular patterns laying near the boundary.
	\item Adaptation of malicious agents to make their outcomes appear like normal in 
	given feature space.
	%TODO \footnote{Adaptation of the agents may }
	\item Normal behavior is evolving in time and thus an normal model defined in one time span 
	can be inaccurate or invalid in future.
	\item An ammount of labeled data needed for derivation of the normal model is insuficient.
	\item Presence of the noise that can  be similar as anomalies, and thus it 
	can be difficult to suppress. 
	\item Different application domains have different notion of an anomaly, 
	thus development of domain-indepedent method is complicated.
\end{itemize}

In general the anomaly detection problem is difficult to solve. 
Most techniques solve a specific formulations of the problem, induced by a 
factors specific for a particular domain. The anomaly detection techniques itself
were developed by adoption of the concepts from diverse disciplines such as \emph{statistics}, 
\emph{machine learning}, \emph{data mining}, \emph{information theory}, \emph{spectral theory}.

\section{Input Data}

Input is generally a collection of data instances, referred as \emph{pattern}, 
\emph{sample} or \emph{observation}.
Each data instance is represented by non-empty set of attributes, also refered as 
\emph{variable} or \emph{feature}.
Attributes can be instances of different data types e.g. \emph{continous}, 
\emph{cathegorical}, or \emph{binary}.
Furthermore in case of each data instance consist of single attribute it is reffered to as
\emph{univariate} otherwise it is \emph{multivariate}. 
For multivariate instances the data types of the attributes might be mixed as well as 
the domain of definition might be different.

\subsection{Relationship Among Data Instances}

Based on presence of the relationship in data, the input data can be further categorized as
\emph{point data}, \emph{sequence data}, \emph{spatial data}, and \emph{graph data}. 
In point data no relationship is assumed among the instances.
In sequence data, presence of the \emph{total order relation}%
\footnote{%
 	In set theory a \emph{total order} is a binary relation on some set $X$. 
	The relation of total order is defined by axioms of \emph{antisymetry}, \emph{transitivity} 
	and \emph{totatlity}. Total order is usually denoted as $\le$.
} %
among data instances is assumed. The sequence data can be time-series, protein sequences, etc.
In \emph{spatial data} presence of \emph{metric}%
\footnote{
	Metric, or distance function, is a non-negative function which defines distance or 
	similarity between elements of the set. Metric is required to satisfy axioms of
	\emph{coincidence}, \emph{symmetry} and \emph{triangle inequality}.
	A metric space is mathematical structure $(X,d)$, where $X$ is a set and function  
	$d:X \times X \rightarrow \mathbf{R}$ is a metric. 
} %
is required. 
The metric determines an neighbourhood of each data instance. The examples of metrics are
\emph{Minkowski metric}% 
\footnote{
	Minkowski metric, defined as $ d(x,y) = (\sum_{i=1}^n(x_i-y_i)^k )^\frac{1}{k}$, 
	is a distance between $n$-vectors $x$ and $y$.
	By choosing value of parameter $k=1$ we get a Mahattan or  a Hamming distance, 
	for $k=2$ we get an Euclid distance, or for $k=\infty $ we get a Chebyshev distance.
} %
(e.g. \emph{Euclidean} distance or \emph{Manhattan} distance), \emph{Levenshtein distance}
(editation distance between strings of characters). Typical example of spatial data is the
coordinate in geographic coordinate system or, asuming our definition, also textual data 
(notice that Levenshtein distance is \emph{metric} among the strings of characters).
The \emph{graph data} instances are represented by graph structure%
\footnote{
	In most common sense, a \emph{graph} $G$ is mathematical structure $G=\left(V,E\right)$
	comprising a set of vertices $V$ with set of edges $E$.
	Edges can be two-element subsets of  $V$ (undirected graph) or ordered pairs 
	of elements of $V$ (directed graph).
	In addition if \emph{weight function} -- $w:E\rightarrow \mathbf{R} $ is defined, 
	assigning a number (e.g. weight, price, etc.) to each edge, we call structure
	$G=\left( V,E,w \right)$ a \emph{weighted graph}.
}%
. As an example of the graph data is map of social social interactions.

In case context are mixed we refer to spatio-temporal (e.g. climate data) or 
graph-temporal data (computer network packet flows).

\subsection{Data Labels}

Labels associated with particular data instances denote if instance is \emph{anomalus}
or \emph{normal}. Labeling is often done by human expert hence it is very expensive and requires
huge effort. Obtaining labels for all possible normal behavior is often less difficult
than obtaining labels for anomalous behavior. Moreover, anomalous behavior is dynamic so
new types of the anomalies might originate. Newly formed anomalies might be then missing 
from models and hence might elude undetected in detection process. 

\section{Anomalies}

Based on presence of the relationship between data instances and problem formulation, anomalies
can be divided into \emph{point anomalies}, \emph{contextual anomalies} and \emph{collective anomalies}.

\subsection{Point anomalies.} In the simplest case, if an individual data instance is considered as
anomalous with respect to the rest of data. No information about relationship between data instances
is assumed. This type of anomaly is target of most of the research studies.
%TODO include examples and pictures

\subsection{Contextual anomalies.} In many cases, an context is present in data set. 
Context is induced by the structure of the data. In case a data instance is anomalous only whithin
a given context, it is called \emph{contextual anomaly}. 
The notion of the context has to be specified within problem formulation. 
By introducing the context in the data features are divided to \emph{contextual features}
and \emph{behavioral features}.

The \emph{contexutal features} are used to determine the context for 
particular data instance. As an examples of the contextual features are: a timestamp denoting 
temporal context in sequential data, a geographic coordinate denoting spatial context.

The \emph{behavioral features} define non-contextual characteristics of an instance. For example, 
the number of arrived packets during network communication whithin a specific time span is considered 
as an behavioral attribute. Identical data instances (in terms of behavioral attributes) 
may be considered as anomalous or non-anomalous in a different contexts.

%TODO some graphics

\subsection{Collective anomalies.} If a collection of related data instances is anomalous with respect
entire dataset it is called \emph{collective anomaly}. The collective anomaly is defined only 
in data set where an relationship among instances are related, e.g. in sequence data, graph data or
spatial data.

\paragraph*{}
It is important to note that \emph{point anomalies} can occur in any data set, while 
\emph{contextual anomalies} depend on notion of the context and its definition in problem formulation, 
and \emph{collective anomalies} are relevant for data where relationship among instances is defined
(e.g. distance metric). So by taking in account the context information a point or collective anomaly
detection problem can be converted into contextual anomaly detection problem.
%TODO example needed 

\section{Techniques}

Typically, the outputs produced by anomaly detection techniques are one of the following two types:
\begin{itemize}
	\item \textbf{Scores.} Assigning score to each data instance depending on the degree to which 
	that instance is considered as anomaly.
	\item \textbf{Labels.} Techniques in this category assign a label (normal or anomalous) to 
	each test instance.
\end{itemize}



Extent in which labels are present in data, can significantly affect mode of operation 
anomaly detection technique:

\begin{itemize}
	\item \textbf{Supervised anomaly detection} sssumes availability of labeled data, for normal
	and also for anomaly class. % imbalanced class distribution
	\item \textbf{Semi-Supervised anomaly detection} assumes availability of data labels only for
	normal class. % I many cases it is not easy to model anomalies
	\item \textbf{Unsupervised anomaly detection} do not require training data. The unsupervised 
	techniques are based on assumption that normal data are more frequent in data than anomalous.	
\end{itemize}

\section{Application domain}



\chapter{Computer Network Security Aspects}

\section{Techniques}

\section{Input Data}

\section{Problem Sets}



\chapter{Proposed Method}

\emph{He et al.} \cite{he2004spectral} showed that the different layers of the network
protocols inprint distinct patterns in a frequency spectrum. 
Further work of \emph{Chen} and \emph{Hwang} \cite{chen2007spectral}
used the features derived %TODO use better word
from frequency spectrum to classify malicious and normal traffic.
They noticed that transport protocols (transmission control protocol --
TCP and user datagram protocol - UDP) 
have distinct frequency spectrum and pover spectral density. 
They exploited this property to identify low-rate attacks on TCP protocol.
%an reduction of service (RoS) attacks.
%TODO more writing about features used
\emph{Wright et al.} \cite{wright2006inferring} researched methods based on 
hidden Markov models to classify different application protocols embeded in 
encrypted application layer.
They developed classification method, able to classify different 
application protocols multiplexed in single encrypted packet flow.
\emph{Dusi et al.} \cite{dusi2009tunnel} brought an statistical approach 
to detect an tunnel inside application layer.
In the paper they described different tunneling techniques and designed 
statistical pattern recognition classifier to identify them.

Our goal is to involve statistical analysis of the frequency components of a 
time-domain signal (spectral analysis) to detection of the tunneled 
connection in application layer.
Our idea is that legitimate connection, that is not missusing application protocol, 
imprints specific patterns in the 
power spectral density that are distinct from spectral density 
of the other unwanted protocols tunneled through application layer. 

In experiments we focused on the application protocols that are most 
spread in computer networks and that are thus most widely miss-used. 
In particular we will analyze protocol HTTP and its encrypted version -- HTTPS. 
Most of the network gateways allow only usage of network protocols HTTP or HTTPS. 
As protocol HTTPS uses encryption, the gateway is unable to detect miss-use by 
analyzing content and thus it is unable to enforce policy restrictions.

Even though it is not possible to analyze payload of particular packets in 
ecrypted connection, it is possible to observe the time of the packet transit, 
its size, direction, source and destination endpoint, etc. 
This data is denoted as packet traces and it is extracted from unencrypted 
part of the packet. 
The goal is to develop feature creation and pattern recognition method for 
the network packet traces involving spectral analysis.
The method is suppoded to infer different application protocols of 
the network traffic only by observation of the packet traces.

\section{Data Collection}

The input data for our method consits of timestamped packet traces. 
Following information has been used from the captured packet data:
\begin{itemize}
	\item packet \emph{timestamp}, i.e. the time of the pass trought the capturing gateway,
	\item identification of source and destination endpoint 
	(i.e. \emph{source address } and \emph{port},
	\emph{destination address}  and \emph{port}
	\footnote{in literature 4-tuple of the  \emph{source address} and 
	\emph{port} and  \emph{destination address} and \emph{port} is 
	denoted as \emph{quad} %TODO cite
	}) and tranmission \emph{protocol} -- we use this infomation (5-tuple) 
	as a key to identify packet \emph{flow},
	\item \emph{size} of packet`s payload,
	\item packet \emph{direction} \footnote{we embed \emph{direction information} 
	into the \emph{size} parameter using negative size  
	if packet travels from destination to source and otherwise positive},
\end{itemize}

\subsection{Training Dataset}

%TODO move to section 5
A \texttt{tcpdump} software \cite{jacobson2009tcpdump}  has been used to 
capture packet data in experimental testbed. 
Testbed consist of three or more peers \emph{A}, \emph{B}, \emph{C}, \emph{D}  etc. 
Peer \emph{A} (a client) is connected with each other using secure tunnel 
to the destination TCP port 443 (destination is on the peers \emph{B},  \emph{C}, ...).
Peer \emph{B} serves virtual private network service and provides network address
translation from the tunnel to the internet.
Peers \emph{C}, \emph{D}, etc. are able to serve single service on the end of tunnel 
(such as HTTP server, telnet, file tranfer service, VOIP, ...).
The traffic tunneled trough peer \emph{B} consist varying number of connections 
and different protocols,
while the others are tunneling single application protocol.

The peer \emph{A} generates traffic to other peers simultaneously, 
mimicking typical behavior of the client of particular application protocol 
(i.e. web browser).
%In addition connections to HTTPS servers are performed from peer \emph{A} directly.
Packet capturing is set up on peer \emph{A} intercepting and storing every packet.
Captured packet flows are then labeled using name of known the application service. 

\subsection{Evaluation Dataset%
\footnote{Evaluation dataset is used during 
assesment of the method on real data.
As we have some old packet trace data, we want to somehow compare it with training samples.
My idea is to prove that exctracted features are epoch invariant by means
statistical goodnes-of-fit test.
On the other hand older protocols can behave in diferent way and thus the 
normal models derived from current data would be
underestimated during assesment proces.
}%
}
The data used for evaluation of proposed method, consists of public packet
traces obtained at....\emph{(to be discussed)}
mixed with simulated and anotated packet traces captured in our testbed.
Before mixing, statistical test has been performed.
The test procedure first extracts and tansforms features
(denoted here as $\vec{f}$) by using our proposed method 
and then compares distribution of features with respect to 
the label -- $Pr\left[ \vec{f} | label \right] $.

Labeling of the obtained dataset is obtained by inferring application 
protocol using destination port.
Since the public dataset can contain malicious connections, 
the result of this statistical test will underestimate the real fitness.
%TODO why?

\section{Feature Creation and Pattern Definition}
%\subsection{Feature exctraction}
For a specified flow $f$ the packet arrival process $x_f\left[t\right]$ 
(or simply \emph{packet process})  is defined as a count of packet arrivals 
at given timespan $I = \left\langle \frac{t}{s}, \frac{t+1}{s} \right)$:
\begin{equation}\label{packetprocess}
	 \forall t \in \mathcal{N} : x_f\left[t\right] = \left| 
	\left\lbrace p : f = flow(p) \wedge timestamp(p) \in I \right\rbrace \right|\, .
\end{equation}
where $s$ is the sample rate, function $flow(p)$  yields the \emph{flow} 5-tuple 
and function $timestamp(p)$  yields the \emph{timestamp} of given packet $p$. 
\footnote{
This is very important formula, we need to focus on it; according 
to \emph{Dusi et al.} \cite{dusi2009tunnel}
zero-length packets are ulikely to be induced by application, 
so we can exclude them here; in addition
they extract incomming and outgoing  stream separately -- 
I think it is good idea to work with in- and out- streams separately 
and compute cross-correlation function instead of auto-correlation function.
The idea of usage I/O cross-correlation function $R_{io}\left(\tau\right)$ 
instead of auto-correlation is that $R_{io}\left(\tau\right)$ (at the specific 
time-lag $\tau$) would enforce the typical request-response round-trip. 
Questionable is, how to physically interpret the resulting spectral components and if the Wiener-Kitchin theorem is applicable.
}

According to a Wiener–Khinchin %TODO reference
theorem the power spectral density $S_{xx}(f)$ is obtained by application of discrete-time 
Fourier transform $\mathcal{F}_{\cdot}(f)$ on autocorrelation function 
of the packet process $R_{xx}\left[\tau\right]$:

\begin{equation}\label{eq:corr}
R_{xx}\left[\tau\right] = E[x\left[t\right]x\left[t+\tau\right]]\, , 
\end{equation}

\begin{equation}\label{eq:psd}
\begin{split}
S_{xx}(f) = \mathcal{F}_{R_{xx}}\left(f\right) = \sum_{\tau=-\infty}^{\infty} 
\left( R_{xx}\left[\tau\right] \exp\left( -\imath 2\pi f\tau \right)\right) \\ 
\forall f \in \left\langle -\frac{s}{2},\frac{s}{2} \right\rangle\, , 
\end{split}
\end{equation}

where $\tau$ is the time-lag, $E\left[\cdot\right]$ is expected value of a random variable and 
$f$ is the frequency. The autocorrelation function is capable of enforcing periodicity.
Equations (\ref{eq:corr}) and (\ref{eq:psd}) hold under assumption that packet process is
\emph{wide-sense stationary random proces}.

As the last assumption seems to be false for infinite time span in network  traffic, 
we involved an \emph{windowing function}  $w(n)$ and 
\emph{discrete Fourier tranform} instead of discrete-time 
Fourier transform. Rectangular windowing is defined as follows:
\begin{equation}
w(n) = \left\lbrace \begin{array}{l} 
1, \mbox{ if } n\in \left\langle 0, M \right) \\ 
0, \mbox{ otherwise} \end{array}\right. \,.
\end{equation}
Windowing function is nonzero inside specified interval $\left\langle 0, M \right)$ 
otherwise it is zero. 
Parameter  $M$ is length of sub-sequence selected from packet arrival proces. 
We iteratively apply windowing function to the whole sequence generating 
non-overlapping adjancent sequence of windows. 
We identify the particular window in this sequence with upper index -- 
e.g. $S_{xx}^i$ is a power spectral density function of an $i$-th window.
Thus, we rewrite equations (\ref{eq:corr}) and (\ref{eq:psd}) for non-overlapping 
windows as follows:
\begin{equation}\label{eq:corr2}
R_{xx}^i\left[m\right] = \frac{1}{M} \sum_{t=0}^{M} x\left[t+iM\right]x\left[t+m+iM\right] \, , 
\end{equation}
\begin{equation}\label{eq:psd2}
\begin{split}
S_{xx}^i(k) = \mathcal{F}_{R_{xx}^i}\left(k\right) = \sum_{m=0}^{M-1}
\left( R_{xx}^i \left[m\right] w(m) \exp\left( -\imath 2\pi m\frac{k}{M} \right)\right)\\
\forall k \in \left\{ 0,1,2,...,M-1 \right\}\, . 
\end{split}
\end{equation}
Note that the windowing function is inherent in equations (\ref{eq:corr2}) and (\ref{eq:psd2})
by using limitted ranges of sumation, and domain of definition of the power spectral density.

By involving windowing function we introduced
\emph{spectral leakge}. Use of diferent windowing function, e.g. \emph{Hann} 
(see equation \ref{eq:hann}),
%or \emph{Hamming} (\ref{eq:hamm}) 
could be appropiate in decreasing leakage.
\begin{equation}\label{eq:hann}
w(n) = \left\lbrace \begin{array}{l} 
0.5\left(1 - \cos \left ( \frac{2 \pi n}{N-1} \right) \right), \mbox{ if } n\in \left\langle 
0, M \right) \\ 0, \mbox{ otherwise} \end{array}\right. \,.
\end{equation}
%\begin{equation}\label{eq:hamm}
%w(n) = \left\lbrace \begin{array}{l} 
%0.54 - 0.46\cos \left ( \frac{2\pi n}{N-1} \right), \mbox{ if } n\in \left\langle 0, M \right) \\
% 0, \mbox{ otherwise} \end{array}\right. \,.
%\end{equation}
Furthermore if the parameter $M$ is too high the packet proces is unlikely to be stationary, 
on the other hand selecting too low value causes that the spectrum is sensitive
to transient phenomena on the network. %TODO why? 

The sampling rate $s$ must be selected according to the Nyquist theorem. 
Too low value entails aliasing%
\footnote{The aliasing is caused by folding of the frequencies above Nyquist frequency
$\frac{s}{2}$ symmetrically below this frequency. Thus this two frequencies are undistinguishable.
To properly reconstruct the signal that contain no frequency higher than $f_{max}$ 
the sample rate is bounded by $s > 2f_{max}$.}%
, while too high value incurs data storage and processing overhead. 

%\subsection{Pattern definition}
By the application of the Fourier transform original features has been mapped into new space.
We denote features in new space as \emph{spectral components} (or \emph{frequency components}).
Temporal context of the features has been altered. 
In new feature space the temporal context is determined by sequence
of the detection windows of size $M$. 
At the level of the detection window a temporal notion is decomposed 
into tuple of temporal functions parametrised by \emph{spectral components}.
Although temporal aspect is still present it has not been used in further analysis.
Anomalies in new feature space are thus regarded to as \emph{point anomalies}. %TODO ?
%The rationale is that the new features decomppsed desired temporal aspects 
%of the original sequential data into declarative form.

There are few parameters affecting quality of the  features: %TODO what?
the smaple rate $s$ and the window length $M$. 
In addition, \emph{spectral components} can be subject to further feature extraction 
or construction, comprising combination of extisting features or discarding irrelevant, 
redundant  or noisy%
\footnote{In case of low signal-to-noise ratio, a feature is typically not usefull 
for discriminative outcome.}%
ones based on the domain knowledge (e.g. sum of a lower resp. a upper half of the 
spectral components resulting in two features a low- and a high-frequency power densities;
or retaining specific spectral features known for its relation to 
periodic stochatic processes to be in research interest).

This aspects and process of seeking of proper parameters are subject of further 
research and they are discussed in \emph{chapter \ref{sec:experiments}}.

\section{Pattern Recognition}

\section{Assessment and Interpretation of Results}


\chapter{Experiments}\label{sec:experiments}

\chapter{Conclusion}



\bibliographystyle{plain}
\bibliography{references}
~\\
~\\

\addcontentsline{toc}{chapter}{Bibliography} 

\cleardoublepage{}

\appendix
\pagenumbering{Roman}
\thispagestyle{empty}  
\renewcommand{\appendixname}{Appendix}%%přílohy, číslování římskými


\chapter*{Appendix}
\addcontentsline{toc}{chapter}{Appendix}

\section*{Whatever}



\end{document}
